\section{Limitations}
\label{sec:conclusion_limitations}

The first limitation is the restricted environment scope.
The experiments used five vector-observation control benchmarks and did not include image observations, language-conditioned environments, multi-agent settings, offline reinforcement learning, model-based planning, or real-world robotic deployment.
The conclusions therefore apply most directly to online, on-policy, simulated control tasks with vector observations.
Tasks with richer sensory input may require different representation-learning assumptions, larger auxiliary models, or different mechanisms for estimating local progress.

The second limitation is that the method was evaluated with PPO as the reinforcement-learning backbone.
PPO provides a stable on-policy setting for controlled comparison, but the behavior of the intrinsic reward may differ under off-policy algorithms, replay buffers, distributed actors, or value-based methods.
The local progress statistics in GLPE are computed from on-policy experience, and their interpretation may change if transitions are repeatedly sampled from a replay buffer or collected by many asynchronous policies.

The third limitation concerns reward shaping.
The augmented reward used during training is not a potential-based transformation of the environment reward and is not guaranteed to preserve the optimal policy of the original task.
Although evaluation uses extrinsic return only, the intrinsic component can still change the policy optimization trajectory and may favor behaviors that are not aligned with the final task objective.
The intrinsic coefficient schedule and clipping range reduce this risk but do not eliminate it.

The fourth limitation is sensitivity to representation quality.
Feature-space impact assumes that distances in the learned latent representation correspond to meaningful transition changes.
Learning progress assumes that forward-model prediction error reflects local competence.
If the encoder fails to preserve action-relevant structure, if the inverse model is insufficiently informative, or if the latent partition mixes unrelated dynamics, then both impact and progress can become unreliable.
This limitation is especially relevant for high-dimensional environments where many different behaviors may map to nearby latent regions.

The fifth limitation is the fixed design of component weights and gate thresholds.
The experiments used a fixed impact weight and environment-specific settings for the learning-progress weight and other hyperparameters.
The ablation results show that different environments favor different components, suggesting that static weights cannot be optimal for all tasks.
Similarly, the gate relies on median-based references, persistence rules, and hysteresis parameters that may not adapt well to every reward structure or dynamics distribution.

The sixth limitation is the computational cost of intrinsic-reward computation.
The method requires an auxiliary encoder, forward model, inverse model, online region statistics, component normalization, and optional gating.
The overhead is especially visible in inexpensive environments where environment interaction is fast.
The gated variant adds further cost through robust threshold computation and gate-state updates.
Although caching can improve throughput, it introduces stale references and can alter gate decisions.

The seventh limitation is the number of random seeds and benchmark scale.
The study reports multiple seeds and uncertainty-aware summaries, but high-variance environments such as \textsc{Humanoid-v5} still exhibit wide confidence intervals.
A larger seed count or a broader benchmark suite would provide stronger evidence about method rankings in difficult continuous-control settings.
The current results are sufficient to identify qualitative behavior and tradeoffs, but they do not prove that small differences between methods are stable across all possible training configurations.

The final limitation is that the study analyzes unproductive curiosity through prediction-error and learning-progress statistics rather than through a formal guarantee.
The gate is designed to suppress regions that empirically appear high-error and low-progress, but this classification can be incorrect.
A region may temporarily show low progress while still being necessary for later task success, or it may appear learnable under the model while remaining irrelevant to extrinsic return.
Thus, the method reduces a known failure mode of curiosity-driven exploration but does not fully solve the broader problem of aligning intrinsic objectives with task objectives.